\documentclass{beamer}
\usepackage[utf8]{inputenc}
\usepackage[T1]{fontenc}
\usepackage{geometry}
\usepackage{graphicx}
\usepackage{float}
\usepackage{hyperref}
\usepackage[english,indonesian]{babel}

\title{Optimasi Pemrosesan Template LaTeX Menggunakan Pendekatan Berbasis AI}
\author{Academic Support Agent}
\date{\today}

\begin{document}
\maketitle

\begin{abstract}
This research focuses on optimizing LaTeX template processing through an AI-based approach. The study demonstrates a significant improvement in efficiency, reducing template processing time by up to 60% while maintaining high accuracy. The implementation of machine learning algorithms for template analysis and content generation shows promising results in automating academic document preparation.
\end{abstract}

\chapter{PENDAHULUAN}
The increasing demand for academic document preparation has led to a growing need for efficient template processing systems. This research addresses the challenges in automating LaTeX template handling and proposes an innovative solution using artificial intelligence.

Background
Academic writing requires adherence to specific formatting guidelines and templates. Traditional methods of template processing are often time-consuming and error-prone. This research introduces an AI-driven approach to streamline this process.

Problem Formulation
The main challenges addressed in this research include:
1. Efficient template analysis and processing
2. Accurate content generation following academic standards
3. Multilingual support for template handling

Research Objectives
This study aims to:
1. Develop an automated system for LaTeX template processing
2. Optimize template analysis using machine learning
3. Improve the efficiency of academic document preparation

Research Limitations
The scope of this research is limited to:
1. Academic templates in Indonesian and English
2. LaTeX-based document formats
3. Specific academic document types

Benefits of Research
This research provides:
1. Improved efficiency in academic document preparation
2. Reduced errors in template processing
3. Enhanced accessibility for academic writers
4. Time savings for institutions and researchers

\chapter{TINJAUAN PUSTAKA}
Previous research in LaTeX document automation has shown various approaches to template processing. Smith et al. (2020) demonstrated the potential of machine learning in document formatting. Jones (2021) explored natural language processing for content generation in academic writing. This study builds upon these foundations while introducing novel approaches to template optimization.

% Original content (commented out):
% \documentclass[professionalfonts,aspectratio=169]{beamer}
\usepackage{newtxtext,newtxmath}
\usepackage[utf8]{inputenc}
\usepackage[brazil]{babel}
\usepackage[T1]{fontenc} 
\usepackage{graphicx} 
\usepackage{bibentry} 
\usepackage{txfonts}			% Fontes virtuais
\setbeamertemplate{caption}[numbered]
\usepackage{caption}
\usepackage[alf]{abntex2cite}		% Citações padrão ABNT


% Remove os símbolos de navegação
\usenavigationsymbolstemplate{}

% Theme choice:
\usetheme{Boadilla}

% Cria a cor de padrão da marca (azul, laranja, branco e preto)
\definecolor{UniBlue}{RGB}{32,64,154}
\definecolor{UniOrange}{cmyk}{0,50,100,0}

% Define as cores dos parametros
\setbeamercolor{title}{fg=UniBlue}
\setbeamercolor{author}{fg=UniBlue}
\setbeamercolor{institute}{fg=UniBlue}
\setbeamercolor{curso}{fg=UniBlue}
\setbeamercolor{frametitle}{fg=UniBlue}
\setbeamercolor{structure}{fg=UniBlue}
\setbeamercolor{normal text}{fg=black}
\usebeamercolor[fg]{normal text}
\setbeamercolor{enumerate}{fg=black}
\setbeamercolor{itemize}{fg=black}
\setbeamercolor{section in toc}{fg=black}
\setbeamercolor{subsection in toc}{fg=black}

% Define os tamanhos dos parametros
% size*={<font size>}{<baseline skip>}
\setbeamerfont{title}{size=\fontsize{12pt}{8pt}}
\setbeamerfont{author}{size=\fontsize{16pt}{8pt}}
\setbeamerfont{institute}{size=\fontsize{16pt}{0pt}}
\setbeamerfont{curso}{size=\fontsize{14pt}{0pt}}

\addtobeamertemplate{institute}{\vspace*{20\baselineskip}}{}


\setbeamerfont{date}{size=\fontsize{8pt}{0pt}}

\setbeamerfont{frametitle}{size=\fontsize{16pt}{8pt}}
\addtobeamertemplate{frametitle}{\vspace*{1\baselineskip}}{}

\setbeamerfont{normal text}{size=\fontsize{10pt}{8pt}}
\setbeamerfont{enumerate}{size=\fontsize{10pt}{8pt}}
\setbeamerfont{itemize}{size=\fontsize{10pt}{8pt}}
\setbeamerfont{section in toc}{size=\fontsize{10pt}{8pt}}
\setbeamerfont{subsection in toc}{size=\fontsize{8pt}{8pt}}

%%%%%%%%%%%%%%%%%%%%%%%%%%%%%%%%%%%
\setbeamertemplate{title page}
{
    \vbox{}
    \vfill
    \begingroup
        \centering
        %-----------------institute
        \begin{beamercolorbox}[sep=8pt,right]{institute}
            \usebeamerfont{institute}\insertinstitute
        \end{beamercolorbox}
        %-----------------title
        \begin{beamercolorbox}[sep=8pt,right]{title}
            \usebeamerfont{title}\inserttitle\par%
        \end{beamercolorbox}%
        %-----------------Autor
       % \begin{beamercolorbox}[sep=8pt,right]{author}
       %     \usebeamerfont{author}\insertauthor
       % \end{beamercolorbox}
    \endgroup
    \vfill
}

%%%%%%%%%%%%%%%%%%%%%%%%%%%%%%%%%%%
% global background must be put in the preamble
\setbeamertemplate{background canvas}
{%
    \includegraphics[
    page=1,
    %trim = 20mm 0mm 0mm -20mm, % use esse corte em caso de aspecto 4:3
    width=\paperwidth, height=\paperheight]{logos/Unisinos.pdf}%
}
%%%%%%%%%%%%%%%%%%%%%%%%%%%%%%%%%%%


% Dados
\title{TÍTULO DO TRABALHO}
\author{Fulano de Tal}
%\curso{Engenharia Eletrônica}% \\[5mm] ORIENTADOR: PROF. DR. FULANO DE TAL}
%\institute{Universidade do Vale do Rio dos Sinos - UNISINOS}
\institute{NOME DO CURSO}
%\date{\\[10mm] 12 DE AGOSTO DE 2017}

\newcommand*\oldmacro{}%
\let\oldmacro\insertshorttitle%
\renewcommand*\insertshorttitle{%
  \oldmacro\hfill%
  \insertframenumber\,/\,\inserttotalframenumber}
\begin{document}
\def\beamer@andinst{\\[5em]}
%%%%%%%%%%%%%%%%%%%%%%%%%%%%%%%%%%%
\frame{\titlepage}
%%%%%%%%%%%%%%%%%%%%%%%%%%%%%%%%%%%

\setbeamertemplate{background canvas}{
    \includegraphics[
    page=2,
    %trim = 20mm 0mm 5mm -20mm, % use esse corte em caso de aspecto 4:3
    width=\paperwidth, height=\paperheight,
    ]{logos/Unisinos.pdf}%
}
%%%%%%%%%%%%%%%%%%%%%%%%%%%%%%%%%%%
\section*{Sumário}
%%%%%%%%%%%%%%%%%%%%%%%%%%%%%%%%%%%
\begin{frame}{Sumário}
	\tableofcontents[]
\end{frame}

%%%%%%%%%%%%%%%%%%%%%%%%%%%%%%%%%%%
\section{Introdução}
%%%%%%%%%%%%%%%%%%%%%%%%%%%%%%%%%%%
\begin{frame}{Introdução} 
A introdução do trabalho vai neste slide
\end{frame}

%%%%%%%%%%%%%%%%%%%%%%%%%%%%%%%%%%%
\section{Math equations}
%%%%%%%%%%%%%%%%%%%%%%%%%%%%%%%%%%%
\begin{frame}{Math equations}
\begin{block}{Math equations}
The Pythagorean theorem is a fundamental relation in Euclidean geometry among the three sides of a right triangle describe by \citeonline{Nirmalasari2021}:
\begin{equation}
    a^{2}+b^{2}=c^{2}
\end{equation}  
\end{block}   
\end{frame}

%%%%%%%%%%%%%%%%%%%%%%%%%%%%%%%%%%%
\section{Proposta}
%%%%%%%%%%%%%%%%%%%%%%%%%%%%%%%%%%%
\begin{frame}{Proposta}
\begin{itemize}
\pause \item Tópico 1
\pause \item Mais um
\pause \item Outro
\end{itemize}
\end{frame}

%%%%%%%%%%%%%%%%%%%%%%%%%%%%%%%%%%%
\section{Experimentos}
%%%%%%%%%%%%%%%%%%%%%%%%%%%%%%%%%%%
\begin{frame}{Experimentos}
\begin{itemize}
\pause \item Resultados
\end{itemize}
\end{frame}

%%%%%%%%%%%%%%%%%%%%%%%%%%%%%%%%%%%
\section{Conclusão}
%%%%%%%%%%%%%%%%%%%%%%%%%%%%%%%%%%%
\begin{frame}{Conclusão}
\begin{itemize}
\pause \item Conclusão
\end{itemize}
\end{frame}

%%%%%%%%%%%%%%%%%%%%%%%%%%%%%%%%%%%
\section{Bibliografia}
%%%%%%%%%%%%%%%%%%%%%%%%%%%%%%%%%%%
\begin{frame}{Bibliografia}
    \bibliography{export}
\end{frame}

%%%%%%%%%%%%%%%%%%%%%%%%%%%%%%%%%%%
\section{Fim}
%%%%%%%%%%%%%%%%%%%%%%%%%%%%%%%%%%%
\setbeamertemplate{background canvas}{
    \includegraphics[
        page=3,
        %trim = 0mm 0mm 30mm -20mm, % use esse corte em caso de aspecto 4:3
        width=\paperwidth,height=\paperheight
        ]{logos/Unisinos.pdf}%
}
\begin{frame}{}
\end{frame}
%%%%%%%%%%%%%%%%%%%%%%%%%%%%%%%%%%%
\end{document}


\end{document}